\documentclass[12pt]{article}
\usepackage[utf8]{inputenc}
\usepackage[T1]{fontenc}
\usepackage{hyperref}
\title{Kurdish Sign Language Language Resources for K-12 Curriculum Development}
\author{Compiled by Claude Sonnet 5.5 (claude-sonnet-5-5) for Andr\'as Kornai}
\date{}
\begin{document}
\maketitle

\section{Introduction}
Kurdish Sign Language (ISO 639-3 zhk) is the sign language of deaf Kurds in the Kurdistan Region of Iraq. Wikipedia describes three dialects, tied to the schools for the deaf in Sulaymaniyah, Erbil and Duhok. It is not mutually intelligible with Iraqi Sign Language. It may have begun with the first Kurdish school for the deaf, founded in Sulaymaniyah in 1982, but its origins are not settled. A 2023 report cited in the search results says the language lacks official recognition, educational support and resources, and recommends a comprehensive dictionary. This suggests none exists yet.

\section{Primary Educational Resources}
No curriculum, textbook or sign-language primer was found. One study describes the language as taught in the schools of the H\^iway Institute in Southern Kurdistan since 1980; the search results gave no details of the materials.

\section{Linguistic and Reference Materials}
\begin{itemize}
\item Wikipedia article on Kurdish Sign Language: \url{https://en.wikipedia.org/wiki/Kurdish_Sign_Language}
\item Deaf Worlds KuSL project page: \url{https://www.deafworlds.org/kusl-project/}
\item University of Raparin journal paper comparing the language with ASL on phonology: \url{https://journal.uor.edu.krd/index.php/JUR/article/download/856/383/1519}
\item Study of the language in the context of the H\^iway Institute: \url{https://gcris.artuklu.edu.tr/entities/publication/5735c0cc-cc06-4ae3-97ff-8421e0371b31}
\end{itemize}

\section{Religious and Community Resources}
No religious resources were found. The Deaf Worlds project page above is the only community-related source found.

\section{Digital and Audio Resources}
\begin{itemize}
\item ``The First Parallel Corpora for Kurdish Sign Language'': \url{https://deepai.org/publication/the-first-parallel-corpora-for-kurdish-sign-language}
\item ``Towards Kurdish Text to Sign Translation'': \url{https://paperswithcode.com/paper/towards-kurdish-text-to-sign-translation}
\item Computers (MDPI) paper presenting a KuSL dataset for recognition: \url{https://www.mdpi.com/2073-431X/13/2/37/htm}
\item Improved recognition of Kurdish Sign Language using a modified CNN: \url{https://www.academia.edu/114305735/Improved_Recognition_of_Kurdish_Sign_Language_Using_Modified_CNN}
\end{itemize}
The search results also mention a HamNoSys/SiGML lexicon dataset by Zina Kamal's group, reported as publicly available on GitHub. I did not see its URL, so it is not cited.

\section{Curriculum Development Context}
Little exists for a K-12 curriculum for deaf and hearing learners. A dictionary, a description of the three dialects and teaching materials would all need to be made, ideally led by deaf signers. The corpora and datasets above are aimed at machine recognition and translation, not teaching. They may still be useful sources of lexical data.

\end{document}